% ============================================================================
\chapter{Дискретная и непрерывная квантовая механика (ДКМ и НКМ), электродинамика Максвелла и граница применимости континуума}
\label{ch:mono-qm-maxwell}
% ============================================================================

\section{Алгебраическое происхождение корпускулярно-волнового дуализма и правила Борна}

В стандартной \textbf{непрерывной квантовой механике (НКМ / CQM)} на континуальном $T$-носителе корпускулярно-волновой дуализм и вероятностное правило Борна $P = |\psi|^2$ принимаются как независимые постулаты, а переход от непрерывной унитарной эволюции Шрёдингера к дискретному срабатыванию детектора («коллапс волновой функции») порождает проблему измерения.

На дискретном кольце $\mathbb{Z}_{512}$ это разделение является прямым следствием теоремы о делении с остатком для разности фаз $\delta = \varphi_y - \varphi_x \in (-512, +512)$:
\begin{align}
  \delta &= 512\,q(\delta) + r(\delta), \label{eq:mono-quotient-remainder} \\
  q(\delta) &\in \{-1, 0, +1\}, \qquad r(\delta) \in (-256, +256]. \label{eq:mono-quotient-remainder-bounds}
\end{align}
\begin{itemize}[leftmargin=*,noitemsep]
\item \textbf{Волновой канал (Канал~1):} главный остаток $r(\delta) \in (-256, +256]$ непрерывно интерферирует вдоль 12 рёбер кубооктаэдра через проекции $\mathcal{C}_3(r)$ и $\mathcal{S}_3(r)$.
\item \textbf{Корпускулярный канал (Канал~2):} целое частное переноса ветви $q(\delta) \in \{-1, 0, +1\}$ (аппаратный бит переноса \texttt{carry\_out} 9-битного сумматора) принимает исключительно дискретные целочисленные значения и отвечает за рождение неделимого топологического кванта (вихревого числа обмотки плакеты).
\end{itemize}

\begin{theorem}[Детерминированный вывод правила Борна $P(\mu) = \mu^2 / 512$ из билинейности оператора голономии $\mathcal{G}_3$]
\label{thm:mono-born-rule}
В уравнениях эволюции поля $\mathcal{G}_3$~\eqref{eq:mono-field-zr}--\eqref{eq:mono-field-draw} операторы парной голономии звена $\mu(\mathbf{x})\,\mu(\mathbf{x}-\mathbf{e})\,\mathcal{C}_3(\Delta\varphi)$ и собственного потенциала ячейки $Q_\square(3)\,\mu(\mathbf{x})^2$ однородно билинейны по амплитудам взаимодействующих узлов. Поэтому для когерентного волнового пакета с целочисленной амплитудой $\mu(\mathbf{x}) = \mu(\mathbf{x}-\mathbf{e}) = \mu \in \{0, 1, \dots, 16\}$ ($\mu^2 \le \Phi_{\mathrm{Pauli}}(3) = 256$) нормированный на квант грани $Q_\square(3) = 128$ фазовый вклад одного звена в ячейку детектора тождественно равен билинейному произведению $\Phi_{\mathrm{wave}}(\mu) = \mu \cdot \mu = \mu^2$. Если собственная фоновая фаза детектора $\varphi_{\mathrm{bg}}$ равномерно пробегает все $\phaseRingCount(3) = 512$ состояний кольца $\mathbb{Z}_{512}$, то число фоновых состояний, при которых сумма $\varphi_{\mathrm{bg}} + \Phi_{\mathrm{wave}}(\mu)$ пересекает границу кольца и рождает корпускулярный бит переноса $q\bigl(\varphi_{\mathrm{bg}} + \Phi_{\mathrm{wave}}(\mu)\bigr) = 1$, тождественно равно квадрату амплитуды:
\begin{align}
  \Phi_{\mathrm{wave}}(\mu) &= \frac{Q_\square(3)\,\mu(\mathbf{x})\,\mu(\mathbf{x}-\mathbf{e})}{Q_\square(3)} \;=\; \mu^2, \label{eq:mono-born-bilinear} \\
  N_{\mathrm{clicks}}(\mu) &= \bigl|\bigl\{\varphi_{\mathrm{bg}} \in \{0,\dots,511\} \;\big|\; \varphi_{\mathrm{bg}} + \Phi_{\mathrm{wave}}(\mu) \ge 512\bigr\}\bigr| \;\equiv\; \mu^2 \;=\; |\psi|^2, \label{eq:mono-born-exact} \\
  P(\mu) &= \frac{N_{\mathrm{clicks}}(\mu)}{\phaseRingCount(3)} \;=\; \frac{\mu^2}{512}. \label{eq:mono-born-prob}
\end{align}
\emph{Формальная верификация:} оператор \texttt{born\allowbreak Click\allowbreak CountOp} и теорема \texttt{reversible\_\allowbreak computing\_\allowbreak zero\_\allowbreak landauer\_\allowbreak heat\_\allowbreak and\_\allowbreak paradox\_\allowbreak resolutions\_\allowbreak operator} в модуле \texttt{lean/\allowbreak Mt\allowbreak Carrier/\allowbreak Rever\allowbreak sible\allowbreak Calculus.lean}.
\end{theorem}

\begin{proof}
Из определения~\eqref{eq:mono-field-zr} и~\eqref{eq:mono-field-draw} при синфазном контакте когерентного пакета амплитуды $\mu(\mathbf{x}) = \mu(\mathbf{x}-\mathbf{e}) = \mu$ ($\Delta\varphi = 0,\ \mathcal{C}_3(0) = Q_\square(3) = 128$) парный голономный вклад одного ребра равен $Q_\square(3)\,\mu\cdot\mu = 128\,\mu^2$, что в единицах шага грани $Q_\square(3) = 128$ даёт фазовый сдвиг $\Phi_{\mathrm{wave}}(\mu) = \mu^2$~\eqref{eq:mono-born-bilinear}. При $0 \le \mu^2 \le 256 < 512$ неравенство $\varphi_{\mathrm{bg}} + \mu^2 \ge 512$ на множестве $\varphi_{\mathrm{bg}} \in \{0, 1, \dots, 511\}$ выполняется тогда и только тогда, когда $\varphi_{\mathrm{bg}} \in \{512 - \mu^2, 512 - \mu^2 + 1, \dots, 511\}$. Число элементов этого отрезка равно $511 - (512 - \mu^2) + 1 = \mu^2$, а деление на полное число равновероятных фоновых фаз $\phaseRingCount(3) = 512$ даёт $P(\mu) = \mu^2 / 512$.
\end{proof}

\section{Порог Гейзенберга \texorpdfstring{$\Delta\varphi_{\mathrm{disc}} = 41$}{dphi = 41}, уравнения Максвелла и граница применимости континуальной теории}

Совокупный фазовый бюджет 13-узлового причинного ядра $B_1(3)$ равен сумме полного кольца $\phaseRingCount(3) = 512$ и числа рёберных связей одной грани кубооктаэдра $\Delta\varphi_{\mathrm{seam}}(3) = \lfloor 256 / 12 \rfloor = 21$:
\begin{equation}
  \label{eq:mono-heisenberg-lock}
  \phaseRingCount(3) + \Delta\varphi_{\mathrm{seam}}(3) \;=\; 512 + 21 \;=\; 533 \;=\; 13 \times 41 \;=\; N_\varphi(3)\cdot \Delta\varphi_{\mathrm{disc}}(3).
\end{equation}
Теорема \texttt{heisenberg\_\allowbreak divisibility\_\allowbreak lock\_\allowbreak unique\_\allowbreak d3\_\allowbreak up\_\allowbreak to\_\allowbreak 8} (\texttt{lean/MtCarrier/Spec.lean}) доказывает, что точная делимость $N_\varphi(d) \mid (\phaseRingCount(d) + \Delta\varphi_{\mathrm{seam}}(d))$ выделяет $d = 3$ и фиксирует фундаментальный порог квантования (мёртвую зону Гейзенберга):
\begin{align}
  \Delta\varphi_{\mathrm{disc}}(3) &= \frac{533}{13} \;=\; 41, \label{eq:mono-heisenberg-41} \\
  \hat{n}_E(\Phi) &= \left\lfloor \frac{|\Phi|}{\Delta\varphi_{\mathrm{disc}}(3)} \right\rfloor \in \{0, 1, \dots, 6\}. \label{eq:mono-heisenberg-quanta}
\end{align}

Всякое подпороговое возмущение $|\Phi| \le 40 < 41$ (образующее класс эквивалентности Гейзенберга мощности $|\mathcal{E}_{\mathrm{Heis}}(3)| = 2\cdot 41 - 1 = 81 = 3^4$) проходит по М-слою носителя когерентно и обратимо с \textbf{нулевым излучением тепловых квантов} ($\hat{n}_E(\Phi) = 0$).

Именно формула~\eqref{eq:mono-heisenberg-41} проводит строгую физическую границу применимости между непрерывной квантовой механикой (НКМ) на $T$-носителе и дискретным $M$-носителем:
\begin{enumerate}[leftmargin=*,noitemsep]
\item \textbf{Область согласия с НКМ ($|\Phi| \gg \Delta\varphi_{\mathrm{disc}} = 41$, $N \gg 1$):}
  надпороговые пакеты подчиняются правилу Борна~\eqref{eq:mono-born-exact}, а кограница главных остатков~\eqref{eq:mono-quotient-remainder} на гранях и кубах решётки тождественно удовлетворяет дискретным уравнениям Максвелла--Фарадея и теореме Гаусса об отсутствии магнитных монополей ($d_1 \circ d_0 = 0$, $\sum_{f \in \partial c} \Phi_f = 0$, теоремы \texttt{d1\_\allowbreak d0\_\allowbreak zero}, \texttt{plaquette\_\allowbreak flux\_\allowbreak eq\_\allowbreak 512\_\allowbreak mul\_\allowbreak winding} и \texttt{discrete\_\allowbreak gauss\_\allowbreak cube\_\allowbreak divergence\_\allowbreak zero} в модуле \texttt{lean/MtCarrier/\allowbreak HolonomyMaxwell.lean}).
\item \textbf{Область расхождения (высокочастотная и подпороговая граница применимости континуума):}
  \begin{itemize}[noitemsep]
  \item \emph{Устранение ультрафиолетовой катастрофы:} так как $|\Phi| \le \Phi_{\mathrm{Pauli}}(3) = 256$, число излучённых квантов энергии с одной ячейки за такт строго ограничено числом квадратных граней кубооктаэдра $\hat{n}_E(\Phi) \le \lfloor 256/41 \rfloor = 6 = |\square(3)|$. Интегралы по импульсам не расходятся на верхнем пределе.
  \item \emph{Преодоление теплового предела Ландауэра ($k_{\mathrm{B}}T \ln 2$):} теорема Ландауэра относится исключительно к необратимому стиранию битов в надпороговом диссипативном Т-канале ($\hat{n}_E \ge 1$). Поскольку микроскопический шаг Верле~\eqref{eq:mono-verlet-law} биективен на $B_2(3) = 72\text{ битах}$, любая логическая операция, выполняемая внутри подпорогового коридора $|\Phi| \le 40$, выделяет ровно $\hat{n}_E = 0$ тепловых квантов.
  \item \emph{Точный закон скрытых микросостояний ($64 \times 81 = 72^2 = 5\,184$):} макроскопический Т-прибор не различает $K_P(3) = 64$ подуровня мантиссы Вейля--Швингера и $|\mathcal{E}_{\mathrm{Heis}}(3)| = 81$ состояние мёртвой зоны Гейзенберга, склеивая ровно $64 \times 81 = 5\,184 = 72^2 = (B_2(3))^2$ детерминированных микросостояний в одно макросостояние. Именно эта проекционная слепота внутреннего прибора порождает кажущуюся термодинамическую стрелу времени (разрешение парадокса Лошмидта) и невозможность демона Лапласа внутри решётки (\texttt{laplace\_\allowbreak demon\_\allowbreak impossibility\_\allowbreak and\_\allowbreak geometric\_\allowbreak maxwell\_\allowbreak demon\_\allowbreak operator}).
  \item \emph{Физическая природа электрического тока как эстафеты дискретного фазовращателя $\mathcal{S}_3(\Delta\varphi)$ и полная система законов теоретической электротехники (DC/AC, Кирхгоф, $RLC$-резонанс, выпрямители и трансформаторы, \texttt{thm5}):}
    в гранецентрированной кубической (ГЦК) решётке металлических проводников ($\mathrm{Cu},\ Z=29$; $\mathrm{Ag},\ Z=47$; $\mathrm{Au},\ Z=79$; $\mathrm{Al},\ Z=13=N_\varphi(3)$) первая координационная сфера из $12$ ближайших соседей каждого узла совпадает с кубооктаэдром $O_2(3)$. При разности фаз $\Delta\varphi \in [0, 128]$ между соседними ячейками чётная квадратурная проекция фазовращателя $V(\Delta\varphi) = \mathcal{C}_3(\Delta\varphi) = 128 - \Delta\varphi$ задаёт потенциальное напряжение звена, а нечётная проекция $I(\Delta\varphi) = \mathcal{S}_3(\Delta\varphi) = \Delta\varphi$ --- электрический ток как скорость передачи фазового кручения от ячейки к ячейке со скоростью света в линии (без механического перемещения частиц сквозь решётку, чем разрешается противоречие между скоростью сигнала $\sim c$ и малой дрейфовой скоростью носителей Друде $v_{\mathrm{drift}} \ll c$). При этом баланс грани $V(\Delta\varphi) + I(\Delta\varphi) \equiv 128 = Q_\square(3)$ определяет передаваемую мощность постоянного тока $P(\Delta\varphi) = V(\Delta\varphi)I(\Delta\varphi) = (128 - \Delta\varphi)\Delta\varphi \le 64^2 = 4\,096$ (максимум согласования импедансов при $\Delta\varphi^* = K_P(3) = 64$, \texttt{thm5}): в подпороговом коридоре $|\Delta\varphi| \le 40 < 41$ фазовая эстафета протекает без рождения фононов ($\hat{n}_E(\Delta\varphi) = 0$, сверхпроводящий/баллистический режим), тогда как при надпороговых скачках $|\Delta\varphi| \ge 41$ на дефектах или тепловых неоднородностях решётки излучается $\hat{n}_E(\Delta\varphi) \in \{1,\dots,6\}$ фононных квантов (закон Джоуля--Ленца). В режиме переменного тока ($\varphi(t)\in\mathbb{Z}_{512}$) законы Кирхгофа для напряжений и токов следуют из $d_1\circ d_0 \equiv 0$, $(d_1 r)_f = 512\,n_f$ и $d_2\circ d_1 \equiv 0$; реактивная мощность конденсатора и индуктивности (сдвинутых на четверть кольца $\pm Q_\square = \pm 128$ тиков) ортогональна активной ($\sum_{\varphi=0}^{511}\mathcal{C}_3(\varphi)\mathcal{S}_3(\varphi) \equiv 0$), а в $LC$-резонансе противофазные реактивные напряжения взаимно аннигилируют в каждой точке кольца ($\mathcal{C}_3(\varphi+128) + \mathcal{C}_3(\varphi+384) \equiv 0$); однополупериодный выпрямитель $\mathrm{rect}_{1/2}(x) = (|x|+x)/2$ даёт за период заряд $\sum_{\varphi=0}^{511}\mathrm{rect}_{1/2}(\mathcal{S}_3(\varphi)) = 128^2 = 16\,384$ (среднее $32 = K_P/2$), двухполупериодный мост Гретца --- заряд $\sum_{\varphi=0}^{511}|\mathcal{S}_3(\varphi)| = 2\times 128^2 = 32\,768$ (средний ток $64 = K_P = I_{\max}/2$), а двухфазный квадратурный мостовой выпрямитель обеспечивает $0\%$ пульсаций без конденсатора ($|\mathcal{C}_3(\varphi)| + |\mathcal{S}_3(\varphi)| \equiv 128$); наконец, в трансформаторе с обмотками $N_1 : N_2$ на общем магнитопроводе ($\Phi_{\mathrm{core}} = 512\,n_f$) закон индукции Фарадея ($V_1 = N_1\varepsilon_1,\ V_2 = N_2\varepsilon_1 \implies N_2 V_1 - N_1 V_2 \equiv 0$) и баланс МДС ($I_1 = N_2 i_0,\ I_2 = N_1 i_0 \implies N_1 I_1 - N_2 I_2 \equiv 0$) дают точное сохранение мощности $V_1 I_1 \equiv V_2 I_2$ и квадратичное трансформирование импеданса $Z_{\mathrm{in}} = (N_1/N_2)^2 Z_{\mathrm{load}}$.
  \end{itemize}
\end{enumerate}

\begin{theorem}[Порог Гейзенберга $\Delta\varphi_{\mathrm{disc}}(3) = 41$, дискретные уравнения Максвелла, скорость связи $c_0 = \sqrt{2}\,c$, граница Цирельсона $2\sqrt{2}$ и законы теоретической электротехники]
\label{thm:mono-heisenberg-maxwell-tsirelson-ee}
На $13$-узловом причинном ядре $B_1(3) = \{\mathbf{0}\}\cup O_2(3)$ и фазовом кольце $\mathbb{Z}_{512}$ выполняются следующие соотношения (\texttt{heisenberg\_\allowbreak divisibility\_\allowbreak lock\_\allowbreak unique\_\allowbreak d3\_\allowbreak up\_\allowbreak to\_\allowbreak 8}, \texttt{d1\_\allowbreak d0\_\allowbreak zero}, \texttt{plaquette\_\allowbreak flux\_\allowbreak eq\_\allowbreak 512\_\allowbreak mul\_\allowbreak winding}, \texttt{discrete\_\allowbreak gauss\_\allowbreak cube\_\allowbreak divergence\_\allowbreak zero}, \texttt{bell\_\allowbreak chsh\_\allowbreak tsirelson\_\allowbreak bound\_\allowbreak from\_\allowbreak cuboctahedron\_\allowbreak link\_\allowbreak speed}, \texttt{thm5}):
\begin{enumerate}[label=\textbf{(\arabic*)},noitemsep]
\item Точная делимость $\phaseRingCount(d) + \Delta\varphi_{\mathrm{seam}}(d)$ на $N_\varphi(d) = 1 + 2d(d-1)$ при $2 \le d \le 8$ выполняется исключительно при $d = 3$ ($512 + 21 = 533 = 13\times 41$), фиксируя $\Delta\varphi_{\mathrm{disc}}(3) = 41$, подпороговый коридор нулевого тепловыделения $|\Phi| \le 40$ ($\hat{n}_E(\Phi) = 0$) и ультрафиолетовый потолок $\hat{n}_E(\Phi) \le \lfloor 256/41 \rfloor = 6 = |\square(3)|$.
\item Поскольку каждое ребро кубооктаэдра $\mathbf{e} \in O_2(3)$ соединяет две координатные оси одновременно ($\|\mathbf{e}\|_2 = \sqrt{1^2 + 1^2 + 0^2}\,dh = \sqrt{2}\,dh$), микроскопическая скорость связи вдоль ребра и граница Цирельсона коррелятора Белла--CHSH для синглетной пары на биссектрисе стационарной фазы Френеля $\psi = K_P(3) = 64 \leftrightarrow \zeta_8 = (\sqrt{2}+i\sqrt{2})/2$ равны:
  \begin{align}
    c_0 &= \sqrt{2}\,c, \qquad c_0^2 \;=\; 2\,c^2, \label{eq:mono-link-speed-c0} \\
    \mathcal{S}_{\mathrm{CHSH}} &= 2\,\frac{c_0}{c} \;=\; 2\sqrt{2} \;>\; 2. \label{eq:mono-tsirelson-bound}
  \end{align}
\item В проводниках с ГЦК-координацией $O_2(3)$ ($\mathrm{Al},\ Z=13$; $\mathrm{Cu},\ Z=29$; $\mathrm{Ag},\ Z=47$; $\mathrm{Au},\ Z=79$) напряжение звена $V(\Delta\varphi) = \mathcal{C}_3(\Delta\varphi) = 128 - \Delta\varphi$ и ток фазовой эстафеты $I(\Delta\varphi) = \mathcal{S}_3(\Delta\varphi) = \Delta\varphi$ при $\Delta\varphi \in [0, 128]$ удовлетворяют балансу грани, закону максимальной согласованной мощности при $\Delta\varphi^* = K_P(3) = 64$, ортогональности активной и реактивной мощностей и тождеству беспульсационного двухфазного квадратурного выпрямителя:
  \begin{align}
    V(\Delta\varphi) + I(\Delta\varphi) &= \mathcal{C}_3(\Delta\varphi) + \mathcal{S}_3(\Delta\varphi) \;\equiv\; 128 \;=\; Q_\square(3), \label{eq:mono-ee-vi-balance} \\
    P(\Delta\varphi) &= V(\Delta\varphi)\,I(\Delta\varphi) \;=\; (128 - \Delta\varphi)\,\Delta\varphi \;\le\; 64^2 \;=\; 4\,096, \label{eq:mono-ee-max-power} \\
    \sum_{\varphi=0}^{511} \mathcal{C}_3(\varphi)\,\mathcal{S}_3(\varphi) &= 0, \qquad
    |\mathcal{C}_3(\varphi)| + |\mathcal{S}_3(\varphi)| \;\equiv\; 128. \label{eq:mono-ee-ortho-rectifier}
  \end{align}
\end{enumerate}
\end{theorem}

\begin{proof}
В пункте~(1) прямое вычисление для $d \in \{2,\dots,8\}$ показывает, что $512 + 21 = 533 = 13\times 41$ делится на $N_\varphi(3) = 13$, а деление с остатком $256 = 6\times 41 + 10$ даёт $\lfloor 256/41 \rfloor = 6$. В пункте~(2) длина вектора $\mathbf{e} \in O_2(3)$ с двумя ненулевыми координатами $\pm 1$ равна $\sqrt{2}$, а на биссектрисе $\psi = 64$ проекции $(\mathcal{C}_3(64), \mathcal{S}_3(64)) = (64, 64)$ дают сумму четырёх корреляторов CHSH $4\times (\sqrt{2}/2) = 2\sqrt{2}$. В пункте~(3) тождество $(128 - \Delta\varphi)\Delta\varphi = 64^2 - (\Delta\varphi - 64)^2 \le 4\,096$ и уравнение $L^1$-алмаза $|\mathcal{C}_3(\varphi)| + |\mathcal{S}_3(\varphi)| = 128$ завершают доказательство.
\end{proof}

\section{Нелинейная механика маятника на произвольных углах, четырёхрежимная теория упругости за пределами закона Гука и предел усталостной выносливости металла (\texttt{MechanicsElasticity.lean})}
\label{sec:mono-mechanics-elasticity-fatigue}

В классической механике и сопротивлении материалов задачи о маятнике и пружине традиционно разделяются двумя искусственными линеаризациями: для маятника постулируется «тонкая нерастяжимая нить» и малость угла ($\sin\varphi \approx \varphi$), а для пружины --- линейный закон Гука $F = k\,\Delta x$, не содержащий ни предела прочности, ни внутреннего трения, ни порога усталостной долговечности металла. В модуле \texttt{lean/MtCarrier/MechanicsElasticity.lean} (\texttt{nonlinear\_\allowbreak elasticity\_\allowbreak four\_\allowbreak regimes\_\allowbreak metal\_\allowbreak fatigue\_\allowbreak and\_\allowbreak coupled\_\allowbreak pendulum\_\allowbreak operator}) доказано, что на дискретном носителе $O_2(3)\times\mathbb{Z}_{512}$ обе задачи описываются единым оператором ротатора $(\mathcal{C}_3, \mathcal{S}_3)$ без малых приближений:

\begin{theorem}[Четыре режима упругости пружины, предел неограниченной выносливости $\Delta\varphi_{\mathrm{disc}}(3)=41$, дислокационный скачок Бюргерса на потолке Паули $\Phi_{\mathrm{Pauli}}(3)=256$ и связанный упругий маятник на произвольных углах, \texttt{MechanicsElasticity.lean}]
\label{thm:mono-mechanics-elasticity-four-regimes}
При $d = 3$ для любой межатомной связи или упругого подвеса с фазовой деформацией $\Delta\varphi \in \mathbb{Z}$, возвращающей силой $F_{\mathrm{spring}}(k, \Delta\varphi) = k\,\mathcal{S}_3(\Delta\varphi)$, потенциальной энергией деформации $U_{\mathrm{spring}}(k, \Delta\varphi) = k\bigl(Q_\square(3) - \mathcal{C}_3(\Delta\varphi)\bigr)$ и оператором внутреннего решёточного трения $\Phi_{\mathrm{fric}}(\Delta\varphi) = \operatorname{heisenbergDeadbandOp}(3, \Delta\varphi)$ выполняются следующие четыре режима (\texttt{0 sorryAx}, \texttt{0 Classical.choice}, \texttt{0 Prop premises}):
\begin{enumerate}[label=\textbf{(\arabic*)},leftmargin=*]
\item \textbf{Режим~I --- Идеально упругий коридор Гука и предел неограниченной усталостной выносливости ($0 \le \Delta\varphi \le \Delta\varphi_{\mathrm{disc}}(3) - 1 = 40$):}
  внутри мёртвой зоны Гейзенберга возвращающая сила и потенциальная энергия строго линейны без приближения Тейлора, а внутреннее трение (возбуждение решётки) и перенос ветви тождественно равны нулю:
  \begin{align}
    F_{\mathrm{spring}}(1, \Delta\varphi) &= \Delta\varphi, \qquad
    U_{\mathrm{spring}}(1, \Delta\varphi) \;=\; \Delta\varphi, \qquad
    \Phi_{\mathrm{fric}}(\Delta\varphi) \;\equiv\; 0, \label{eq:mono-hooke-lossless} \\
    D_{\mathrm{fatigue}}(N_{\mathrm{cycles}}, 40) &:= N_{\mathrm{cycles}}\cdot|\Phi_{\mathrm{fric}}(40)| \;\equiv\; 0 \qquad (\forall\,N_{\mathrm{cycles}} \in \mathbb{N}), \label{eq:mono-infinite-fatigue-endurance}
  \end{align}
  что определяет строгий модельный порог неограниченной усталостной выносливости ($\sigma_{-1} \leftrightarrow \Delta\varphi_{\mathrm{disc}} = 41$), ниже которого усталостное повреждение не накапливается при сколь угодно большом числе циклов нагружения.
\item \textbf{Режим~II --- Неупругое внутреннее трение и абсолютный предел прочности на четверти оборота ($41 \le \Delta\varphi \le Q_\square(3) = 128$):}
  при $\Delta\varphi \ge 41$ открывается пороговый вентиль $\Phi_{\mathrm{fric}}(\Delta\varphi) = \Delta\varphi \ne 0$ ($\hat{n}_E(\Delta\varphi) = \lfloor\Delta\varphi/41\rfloor \ge 1$), переводящий часть энергии каждого цикла качания во внутреннее дрожание атомов кристаллической решётки (внутреннее трение и акустическую эмиссию в металле даже в полном вакууме), а в седловой точке $\Delta\varphi = Q_\square(3) = 128$ возвращающая сила достигает глобального максимума по всему кольцу (предела упругой прочности связи):
  \begin{equation}
    \label{eq:mono-uts-peak-128}
    F_{\mathrm{spring}}(1, 128) \;=\; Q_\square(3) \;=\; 128 \;\ge\; F_{\mathrm{spring}}(1, \varphi) \quad (\forall\,\varphi \in \mathbb{Z}_{512}), \qquad
    \mathcal{C}_3(128) \;=\; 0.
  \end{equation}
\item \textbf{Режим~III --- Зона разупрочнения и образования шейки ($128 < \Delta\varphi < \Phi_{\mathrm{Pauli}}(3) = 256$):}
  за перевалом четверти оборота $\mathcal{C}_3(\Delta\varphi) < 0$ (антисвязывающая ветвь), а наклон возвращающей силы меняет знак на отрицательный ($\Delta\mathcal{S}_3 = -1 < 0$), так что при дальнейшем росте удлинения сила сопротивления падает от $F_{\mathrm{spring}}(1, 129) = 127$ до $F_{\mathrm{spring}}(1, 255) = 1$.
\item \textbf{Режим~IV --- Потолок текучести/разрыва Паули ($\Delta\varphi = 256$) и рождение квантованной дислокации Бюргерса ($\Delta\varphi = 257$):}
  на антиподальном потолке Паули $\Delta\varphi = \Phi_{\mathrm{Pauli}}(3) = 256$ возвращающая сила обращается в нуль ($F_{\mathrm{spring}}(1, 256) = 0,\ U_{\mathrm{spring}}(1, 256) = 256$), а переход через границу ветви $\Delta\varphi = 257$ вызывает скачок частного переноса $q_e : 0 \to +1$ и рождение в плакете решётки топологического вихря (дислокации с целочисленным вектором Бюргерса):
  \begin{equation}
    \label{eq:mono-burgers-dislocation-jump}
    q_e(256) \;=\; 0 \;\longrightarrow\; q_e(257) \;=\; +1, \qquad
    w_f(0, 128, 256, 257) \;=\; +1, \qquad
    \mathcal{B}_f \;=\; 512.
  \end{equation}
\item \textbf{Связанный упругий маятник на произвольных углах от $0^\circ$ до $360^\circ$ без приближения $\sin\varphi \approx \varphi$:}
  продольное натяжение упругого подвеса $T_{\mathrm{long}}(g, \varphi) = g\,\mathcal{C}_3(\varphi)$ (растягивающее нить по закону Режимов~I--IV) и поперечная возвращающая сила $F_{\mathrm{trans}}(g, \varphi) = g\,\mathcal{S}_3(\varphi)$ при любом угле отклонения $\varphi \in \mathbb{Z}_{512}$ удовлетворяют точному закону сохранения $L^1$-нормы и постоянству разностного вронскиана:
  \begin{align}
    |T_{\mathrm{long}}(1, \varphi)| + |F_{\mathrm{trans}}(1, \varphi)| &\equiv Q_\square(3) \;=\; 128, \label{eq:mono-pendulum-l1-conservation} \\
    \bigl(T_{\mathrm{long}}(g, 0),\ F_{\mathrm{trans}}(g, 0)\bigr) &= (128g,\ 0), \qquad
    \bigl(T_{\mathrm{long}}(g, 128),\ F_{\mathrm{trans}}(g, 128)\bigr) \;=\; (0,\ 128g), \nonumber \\
    \bigl(T_{\mathrm{long}}(g, 256),\ F_{\mathrm{trans}}(g, 256)\bigr) &= (-128g,\ 0), \qquad
    \mathcal{K}_{\mathrm{ell}}(1, 1) \;=\; 2\,\Phi_{\mathrm{Pauli}}(3)\,Q_\square(3) \;=\; 65\,536. \label{eq:mono-pendulum-jacobi-period}
  \end{align}
\end{enumerate}
\end{theorem}

\begin{proof}
Все утверждения Теоремы~\ref{thm:mono-mechanics-elasticity-four-regimes} формально проверены ядром Lean~4 в теореме \texttt{nonlinear\_elasticity\_four\_regimes\_metal\_fatigue\_and\_coupled\_pendulum\_operator} модуля \texttt{lean/MtCarrier/MechanicsElasticity.lean}.
\end{proof}

\begin{quote}
\small\textbf{Методологическое примечание о границе между формальной моделью и натурными испытаниями материалов.}
Соотношения~\eqref{eq:mono-hooke-lossless}--\eqref{eq:mono-pendulum-jacobi-period} строго доказаны в ядре Lean~4 как точные алгебраические свойства дискретной модели $O_2(3)\times\mathbb{Z}_{512}$. При этом прямой пересчёт модельных порогов одной связи ($\Delta\varphi_{\mathrm{disc}} = 41$, $Q_\square = 128$, $\Phi_{\mathrm{Pauli}} = 256$) в макроскопические инженерные напряжения ($\sigma_{-1}, \sigma_{0{,}2}, \sigma_{\mathrm{B}}$ в МПа) и кривые усталости Вёлера конкретных поликристаллических авиационных и конструкционных сплавов (с учётом границ зёрен, легирования и концентрации напряжений) \textbf{ещё не проходил независимой внешней экспериментальной проверки} и составляет открытую программу стендовых испытаний на вибростендах в вакууме с синхронной регистрацией акустической эмиссии и тепловыделения.
\end{quote}

